\input{../../stock/en-tete_v6.tex}

\newcommand{\titrechapitre}{titre}
\newcommand{\numerochapitre}{numéro}

\renewcommand{\headrulewidth}{0.4pt}
\renewcommand{\footrulewidth}{0.4pt}
\fancyfoot[L]{Notes de Raphaël JONTEF}
\fancyfoot[C]{\thepage}
\fancyfoot[R]{Enseirb-Matmeca}
\fancyhead[L]{Cours de Mathématiques}
\fancyhead[R]{\titrechapitre}

\begin{document}
\input{../../stock/commands.tex}

% Page de titre custom
\setlength{\headheight}{15pt}
\begin{titlepage}
    \centering
    \vspace*{3cm}
    {\huge Chapitre \numerochapitre\par}
    {\Huge \bfseries\titrechapitre\par}
    \vspace{2cm}
    {\Large Notes rédigées par Raphaël JONTEF\par}
{\Large Avec l'aide de \par}
    \vspace{0.5cm}
    {\normalsize Cours enseigné par \par}
    \vspace{4cm}
    {\small Enseirb-Matmeca\par}
    {\small filière informatique\par}
    \vfill
    {\large \today\par}
\end{titlepage}

% plan du cours
\tableofcontents
\newpage


% -----------Corps du document--------------

\section{Résolution de systèmes linéaires}

Soit $A \in \matrice{n}(\R)$, $B \in \R^n$
En pratique, la solution numérique $x + \delta x$ de l'égalité vectorielle $AX = B$ vérifie~:
$$A\times(x+\delta x) = b + \delta b$$
On a alors~:
$$\delta b = A\times(x + \delta x) - b$$
Le vecteur $x$ étant inconnu, on ne peut pas calculer directement $\delta b$. On cherche donc à majorer $\norme{\delta b}$.\\\par

\subsection{Rappels sur les normes usuelles sur $\R^n$}

\begin{definition}{}{norme infinie}
    Pour $X \in \R^n$, on définit la norme infinie par~:
    $$\norme{X}_\infty = \max_{1 \leq i \leq n} |x_i|$$
    En Python, on peut la calculer avec la fonction \texttt{numpy.linalg.norm(X, inf)}.
\end{definition}

\begin{definition}{}{norme de Mannhattan}
    Pour $X \in \R^n$, on définit la norme de Manhattan par~:
    $$\norme{X}_1 = \sum_{i=1}^n |x_i|$$
    En Python, on peut la calculer avec la fonction \texttt{numpy.linalg.norm(X, 1)}.
\end{definition}

\begin{definition}{}{norme euclidienne}
    Pour $X \in \R^n$, on définit la norme euclidienne par~:
    $$\norme{X}_2 = \left(\sum_{i=1}^n |x_i|^2\right)^{1/2}$$
    En Python, on peut la calculer avec la fonction \texttt{numpy.linalg.norm(X)}.
\end{definition}

\begin{definition}{}{norme subordonnée}
    Pour $A \in \matrice{n}(\R)$, on définit la norme subordonnée associée à une norme vectorielle $\norme{\cdot}$ par~:
    $$\norme{A} = \sup_{X \in \R^n, X \neq 0} \frac{\norme{AX}}{\norme{X}}$$
    En Python, on peut la calculer avec la fonction \texttt{numpy.linalg.norm(A, ord)} où \texttt{ord} est l'ordre de la norme vectorielle associée (\code{1}, \code{2} ou \code{inf}).
    
\end{definition}


\begin{remarque}{}{sur la norme euclidienne d'une matrice}
    On a $\norme{A}_2 = \sqrt{\rho(AA^\top)}$ où $\rho(M)$ est la valeur propre de $M$ de plus grande valeur absolue.
    Cette norme est difficile à calculer en pratique. On préfère donc utiliser les normes $\norme{\cdot}_1$ ou $\norme{\cdot}_\infty$. qui sont plus rapides à calculer.
    
\end{remarque}

\subsection{Erreur d'approximation et méthode d'estimation de cette erreur}

Si $A$ est inversible, on a~:
\begin{align*}
    b &= Ax
    &\Rightarrow \norme{b} &= \norme{Ax} \leq \norme{A}\norme{x}\\
\end{align*}

De même~:
\begin{align*}
    A \times (x + \delta x) &= b + \delta b \\
    &\Rightarrow A \delta x = \delta b\\
    &\Rightarrow \norme{\delta b} = \norme{A \delta x} \leq \norme{A}\norme{\delta x}\\
    &\Rightarrow \norme{b} \norme{\delta b} \leq \norme{A}\norme{x} \norme{A^{-1}}\norme{\delta x}\\
    &\Rightarrow \underbrace{\frac{\norme{\delta x}}{\norme{x}}}_{\text{erreur relative sur la solution}} \leq \underbrace{\norme{A}\norme{A^{-1}}}_{\mr{cond}(A)} \underbrace{\frac{\norme{\delta b}}{b}}_\text{erreur que l'on mesure}\\
\end{align*}
L'erreur est amplifiée d'un facteur appelé \notion{conditionnement} de la matrice $A$, noté~:
$$\mr{cond}(A) = \norme{A}\norme{A^{-1}}$$
En Python, on le calcule avec \code{cond(A, ord)} (\textit{cf.} plus haut).

\begin{definition}{}{mauvais conditionnement}
    On dit d'une matrice $A \in \matrice{n}(\R)$ qu'elle est \notion{mal conditionnée} si son conditionnement est élevé~:
    \begin{itemize}
        \item Si $\mr{cond}(A) = 10^4$, alors $A$ est mal conditionnée
        \item Si $\mr{cond}(A) = 3$, aors $A$ est bien conditionnée.
    \end{itemize}
\end{definition}

Moralement, pour résoudre $Ax = b$~:
\begin{itemize}
    \item On trouve $x + \delta x$
    \item On évalue le \notion{résidu} $\norme{\delta b} = \norme{b - A\times(x + \delta x)}$
\end{itemize}

\begin{remarque}{}{mauvais conditionnement}
    Les matrices obtenues en discrétisant des équations physiques sont généralement mal conditionnées.
\end{remarque}

\section{Principes généraux pour réduire les erreurs d'arrondi}

\subsection{Éviter de soustraire deux quantités absoluements proches}

\begin{exemple}{}{soustraction au cosinus}
    POur calculer $y = 1-\cos(x)$, pour $x$ proche de 0, on peut utiliser la formule~:
    $$\cos(2x) = 1 - 2\sin^2(x)$$
    pour obtenir~:
    $$y = 2 \sin^2(x/2)$$
\end{exemple}

\begin{exemple}{}{}
    On s'intéresse à la recherche des racines de $P = aX^2 + bX + c$, de discriminant $\Delta$ vérifiant~:
    $$\Delta \gg 0$$
    On a alors~:
    $$b^2 \gg 4ac$$
    Ce discriminant est donc proche de $b^2$. Ainsi, quand on calcule~:
    $$x_1 = \frac{-b-\Delta}{2a} \qquad \text{et} \qquad x_2 = \frac{-b + \Delta}{2a}$$ 
    le calcul de $x2$ est imprécis à cause du dénominateur. Mais via les formules de Viète, $x_2 = \frac{c}{ax_1}$.
\end{exemple}

\subsection{Sommer d'abord les termes les plus petits}
Sommer un grand nombre avec un petit nombre entraîne la troncature des plus petites décimales du petit nombre.
\begin{exemple}{}{calcul d'une somme}
    Pour calculer~:
    $$S = \sum_{i=1}^n \frac{1}{i}$$
    bien qu'il soit naturel de sommer les termes dans l'ordre croissant des indices, l'ordre rétrograde est ici à privilégier. On calcule sur quatre décimales~:\\
    \centering
    \begin{tabular}{c|c|c|c}
        $n$ & croissant & rétrograde & exact\\
        $100$ & $5,187$ & $5,187$ & $5,1873775$\\
        $1\,000$ & $7,449$ & $7,485$ & $7,485471$\\
        $10\,000$ & $8,449$ & $9,448$ & $9,787606$\\
        
    \end{tabular}
\end{exemple}

\section{Calculs de complexité}
Les additions, soustractions et multiplication ont aujourd'hui des coûts temporels semblables, mais la division est plus complexe. On cherche à s'en passer le plus possible.

\begin{exemple}{}{complexité du produit d'une matrice avec un vecteur}
    Soit $A \in \matrice{n,p}(\R)$ et $x \in \R^n$
    Notons $y = Ax$ et $y = (y_1,\, \dots,\, y_n)^\top$. On a~:
    $$\forall i \in \intint{1}{n},\, y_i = \sum_{j=1}^p a_{i,j} x_j$$.
    Le calcul de chaque $y_i$ demande donc $2p$ opérations, ainsi le calcul de $y$ demande $2np$ opérations.
    
\end{exemple}

\begin{exemple}{}{résoltuion d'un système linéaire triangulaire}
    Calculons la complexité de la résolution de $Tx = b$. On prend $T$ inférieure
\end{exemple}



\end{document}